% ── Cieľ práce ────────────────────────────────────────────────
\chapter{Cieľ práce}

        Cieľom tejto práce bolo navrhnúť a~implementovať webovú aplikáciu na správu projektov
s~integrovaným AI~chatom nad dokumentmi. Výsledný systém má nahradiť roztrieštenú prácu
s~viacerými nástrojmi a~poskytnúť tímom jednotné prostredie pre riadenie úloh, ukladanie
dokumentov a~inteligentnú prácu s~ich obsahom. Aplikáciu sme navrhli tak, aby bola nasaditeľná
lokálne na infraštruktúre organizácie bez závislosti na externých cloudových službách na
spracovanie dokumentov.

Na splnenie hlavného cieľa sme najprv navrhli architektúru systému ako modulárne
riešenie pozostávajúce z~backendu, frontendu, databázovej vrstvy a~AI~modulu, aby
bolo možné funkcionalitu v~budúcnosti jednoducho rozširovať. Následne sme implementovali
správu projektov a~úloh so zameraním na vytváranie projektov, evidenciu úloh, priraďovanie
zodpovedných osôb, sledovanie stavu a~podporu základného tímového workflow. Súčasťou
riešenia je aj správa dokumentov, ktorá umožňuje nahrávanie, ukladanie a~organizáciu
súborov v~kontexte projektu s~podporou formátov PDF a~DOCX. Dôležitou časťou práce
bola integrácia AI~chatu nad dokumentmi, teda sprístupnenie otázok v~prirodzenom jazyku
nad obsahom vybraných dokumentov a~vracanie relevantných odpovedí na základe extrahovaného
textu. Zároveň sme zabezpečili, aby aplikácia fungovala bezpečne v~lokálnej prevádzke,
dokumenty ostávali v~internom prostredí organizácie a~komunikácia prebiehala cez zabezpečený
HTTPS kanál. Posledným krokom bolo overenie funkčnosti riešenia prostredníctvom praktického
testovania hlavných scenárov používania a~zhodnotenia, či systém spĺňa stanovené funkčné
a~nefunkčné požiadavky.

Výstupom práce je funkčný prototyp webovej aplikácie, ktorý demonštruje,
že prepojenie projektového riadenia, dokumentového manažmentu a~AI~asistencie
v~jednom systéme je realizovateľné aj v~podmienkach lokálnej infraštruktúry.
